\subsection{A fixed Stein hypersurface for all Finsler metrics}
\label{subsec:fixed-stein-finsler-hypersurface}
\index{Stein manifold!fixed Finsler test hypersurface}

\Cref{cor:hirzebruch-stable-ample-finsler-obstruction} shows that no strongly pseudoconvex Finsler metric on $E$ has semipositive Kobayashi curvature, but it does not locate the failure. For a single metric, a point where the curvature is not semipositive gives a local hypersurface through that point on which the metric has a nondegenerate critical point with many positive Levi directions. Holomorphic approximation on a Stein manifold allows one hypersurface to serve for all metrics at once.

In this subsection, $X$ is a connected smooth projective variety of dimension $n\geq1$, and $E\to X$ is an algebraic vector bundle of rank $r\geq2$, viewed as a holomorphic bundle. We use the squared Finsler metrics and the convention for Kobayashi curvature of \cref{subsec:polar-finsler-convexity}. A compact subset $K$ of a Stein manifold $Y$ is \emph{holomorphically convex} if it coincides with its holomorphic hull $\{y\in Y:|f(y)|\leq\sup_K|f|\text{ for all }f\in\mathcal O(Y)\}$. We use the classical Oka--Weil theorem in the following form: a function holomorphic near a holomorphically convex compact subset of a Stein manifold can be approximated uniformly there by global holomorphic functions. We also use polynomial Runge approximation in one variable on a union of two disjoint closed disks.

\Needspace{6\baselineskip}
\begin{lemma}\label{lem:separated-holomorphic-convex-compacts}
\index{Holomorphic convexity!separated compact sets}
Let $K,L$ be holomorphically convex compact subsets of a Stein manifold $Y$. If a global holomorphic function maps them into disjoint closed disks in $\mathbb C$, then $K\cup L$ is holomorphically convex.
\end{lemma}
\begin{proof}
Write $h$ for the function and $\Delta_K,\Delta_L$ for the disks. Their union is polynomially convex, so testing with polynomials in $h$ shows that $h$ maps the holomorphic hull of $K\cup L$ into $\Delta_K\cup\Delta_L$. Suppose that $p\notin K\cup L$ and $h(p)\in\Delta_K$. Holomorphic convexity of $K$ gives $a\in\mathcal O(Y)$ with $|a(p)|>\sup_K|a|$. Choose polynomials tending uniformly to one on $\Delta_K$ and to zero on $\Delta_L$. Their compositions with $h$, multiplied by $a$, converge to $a$ at $p$ and on $K$, and to zero on $L$. A sufficiently good approximation therefore separates $p$ from $K\cup L$. The case $h(p)\in\Delta_L$ is identical, with $K$ and $L$ interchanged. Thus no point outside $K\cup L$ belongs to its hull.
\end{proof}

\Needspace{6\baselineskip}
\begin{lemma}\label{lem:dilated-countable-holomorphic-tests}
\index{Holomorphic approximation!countable dilated tests}
Let $U$ be a connected smooth affine variety and let $q\geq0$ be an integer. Put $B=U\times\mathbb C^*\times\mathbb C^q$, with coordinates $(x,t,s)$, and for $\lambda\in\mathbb C^*$ let $d_\lambda(x,t,s)=(x,\lambda t,\lambda s)$. For each $j\geq1$, let $C_j\subset B$ be a holomorphically convex compact set, let $g_j$ be holomorphic near $C_j$, and let $\delta_j>0$. Suppose that $t(C_j)$ lies in the disk $|t-t_j|\leq|t_j|/4$ for some $t_j\neq0$. Then there are $\lambda_j\in\mathbb C^*$ and $f\in\mathcal O(B)$ such that
\[
 \sup_{d_{\lambda_j}(C_j)}
 |f-\lambda_j g_j\circ d_{\lambda_j}^{-1}|<|\lambda_j|\delta_j
 \qquad(j\geq1).
\]
\end{lemma}
\begin{proof}
Choose an exhaustion $(L_j)_{j\geq0}$ of the affine, hence Stein, manifold $B$ by holomorphically convex compact sets with $L_j\subset\operatorname{int}L_{j+1}$. We construct $f_j\in\mathcal O(B)$ and compact sets $K_j$ inductively, starting with $f_0=0$ and $K_0=L_0$.

At stage $j\geq1$, choose $M>0$ with $t(K_{j-1})\subset\{|t|\leq M\}$, and choose $R>4M/3$. Dilation by $\lambda_j=R/t_j$ moves the $t$-projection of $C_j$ into the disk $|t-R|\leq R/4$, which is disjoint from $|t|\leq M$. For the rest of the induction we replace $C_j$, $g_j$ and $\delta_j$ by the dilated compact set $d_{\lambda_j}(C_j)$, the function $\lambda_jg_j\circ d_{\lambda_j}^{-1}$ and the tolerance $|\lambda_j|\delta_j$, keeping the same notation. \Cref{lem:separated-holomorphic-convex-compacts}, applied with the function $t$, makes $K_{j-1}\cup C_j$ holomorphically convex.

On disjoint neighborhoods of $K_{j-1}$ and $C_j$, prescribe $f_{j-1}$ and $g_j$, respectively. Oka--Weil approximation gives $f_j\in\mathcal O(B)$ with
\[
 \sup_{K_{j-1}}|f_j-f_{j-1}|<\epsilon_j,
 \qquad \sup_{C_j}|f_j-g_j|<\epsilon_j,
 \qquad
 \epsilon_j=2^{-j-2}\min(1,\delta_1,\ldots,\delta_j).
\]
Let $K_j$ be a member of the fixed exhaustion containing $K_{j-1}\cup C_j\cup L_j$ in its interior. The interiors of the sets $K_j$ exhaust $B$, and the summable errors make $f_j$ converge locally uniformly to some $f\in\mathcal O(B)$. For a fixed $j$, every later approximation controls its error on $C_j$, since $C_j\subset K_{m-1}$ for $m>j$. Consequently
\[
 \sup_{C_j}|f-g_j|
 \leq\sum_{m\geq j}\epsilon_m
 \leq2^{-j-1}\delta_j<\delta_j.
\]
Returning to the original notation gives the stated tolerances.
\end{proof}

\Needspace{6\baselineskip}
\begin{theorem}\label{thm:fixed-stein-finsler-hypersurface}
\index{Finsler metric!fixed hypersurface test}
\index{Kobayashi curvature!critical hyperplane test}
\index{Critical point!holomorphic quadratic jet}
\index{Levi index!fixed hypersurface test}
Suppose that $E$ admits no smooth strongly pseudoconvex squared Finsler metric with semipositive Kobayashi curvature. Then there is a connected smooth Stein hypersurface $\mathcal S\subset E^\times$, locally closed in the analytic topology, with the following property: for every smooth strongly pseudoconvex squared Finsler metric $G$ on $E$, the restriction $G|_{\mathcal S}$ has a nondegenerate critical point with at least $r$ positive Levi eigenvalues, and hence with at most $n-1$ negative ones.
\end{theorem}
\begin{proof}
\emph{Step 1: coordinates.} An algebraic frame trivializes $E$ over a nonempty affine open subset $U\subset X$. Since $X$ is smooth and connected, it is irreducible, so $U$ is dense and connected. Write the fiber coordinates as $(w,t,s)\in\mathbb C\times\mathbb C\times\mathbb C^{r-2}$, and put $B=U\times\mathbb C^*\times\mathbb C^{r-2}$. We seek $\mathcal S$ as a graph $w=f(x,t,s)$ over $B$; such a graph avoids the zero section because $t\ne0$.

\emph{Step 2: a local test for one metric.} Fix a candidate metric $G$. Since its Kobayashi curvature is not semipositive, its horizontal Schur complement has a positive direction somewhere. This condition is open, so it holds at a point $v_0=(x_0,w_0,t_0,s_0)$ with $x_0\in U$, $t_0\neq0$ and $G_w(v_0)\neq0$. Indeed, the first two conditions define a dense open set, and the zero set of $G_w$ has empty interior, since vanishing of $G_w$ on an open set would give $G_{w\bar w}=0$ there, contradicting the positivity of the vertical Hessian.

Let $H$ be the Levi-orthogonal complement of the vertical tangent space $V$ at $v_0$. The Euler identities in the proof of \cref{lem:polar-finsler-osculating-section} give
\[
 K:=\ker\partial G(v_0)=H\oplus(V\cap K),
 \qquad \dim_{\mathbb C}(V\cap K)=r-1.
\]
Indeed, $\partial G$ vanishes on $H$, and its value on the vertical radial vector is $G(v_0)>0$. The two summands are Levi orthogonal, the vertical summand is positive definite, and the horizontal summand has a positive direction. Thus the Levi form of $G$ restricted to $K$ has at least $r$ positive eigenvalues.

Since $G_w(v_0)\neq0$, the space $K$ is the graph of a linear map defined on $T_{b_0}^{1,0}B$, where $b_0=(x_0,t_0,s_0)$. Choose a local holomorphic function $g$ near $b_0$ with $g(b_0)=w_0$ whose graph has tangent space $K$ at $v_0$. Then $F=G\circ\operatorname{Graph}(g)$ has a critical point at $b_0$, and its Levi form there is the restriction of the Levi form of $G$ to $K$. Keep this first jet of $g$ fixed and vary its quadratic jet. The mixed second derivatives of $F$ do not change, while in local coordinates $z$ on $B$ the pure second derivatives satisfy
\[
 F_{\alpha\beta}(b_0)=A_{\alpha\beta}+G_w(v_0)g_{\alpha\beta}(b_0),
\]
where $A$ is fixed. Since $G_w(v_0)\ne0$, we can prescribe the pure quadratic part of $F$ to be $\tau\operatorname{Re}\sum_\alpha z_\alpha^2$, whose real Hessian has eigenvalues $\pm2\tau$. If $\tau$ is larger than half the operator norm of the real Hessian of the fixed mixed part, the full real Hessian is invertible by the Neumann-series criterion. The critical point is then nondegenerate, with the same Levi form.

\emph{Step 3: stability of the test.} Choose compact neighborhoods $C'\subset\operatorname{int}C$ of $b_0$ on which $g$ is defined, with $C$ holomorphically convex and $t(C)\subset\{|t-t_0|\leq|t_0|/4\}$. For a closed affine embedding of $B$, the intersection of $B$ with a small closed Euclidean ball centered at $b_0$ is such a set $C$, since a closed Euclidean ball is convex and hence polynomially convex. Fix a smooth Hermitian background metric on $E$ and its compact unit sphere bundle $\Sigma$. By homogeneity, the $C^2$ distance between metrics on $\Sigma$ controls their $C^2$ distance on each fixed compact subset of $E^\times$. Persistence of nondegenerate critical points and openness of the conditions on the Levi eigenvalues therefore give a neighborhood $\mathcal V_G$ of $G|_\Sigma$ for the $C^2$ topology and a number $\delta_G>0$ with the following property. For every metric whose restriction to $\Sigma$ lies in $\mathcal V_G$ and every holomorphic function within $\delta_G$ of $g$ uniformly on $C$, the restriction of the metric to the graph of the function still has a nondegenerate critical point over $C'$ with at least $r$ positive Levi eigenvalues. Here the Cauchy estimates on $C'$ turn uniform closeness on $C$ into $C^2$ closeness on $C'$.

\emph{Step 4: one hypersurface for all metrics.} The restrictions of the candidate metrics to $\Sigma$ form a subspace of the separable metrizable space $C^2(\Sigma)$. This subspace is second countable, hence Lindel\"of, so the cover by the neighborhoods $\mathcal V_G$ has a countable subcover $(\mathcal V_j)_{j\geq1}$, with associated data $C_j,g_j,\delta_j$. For $D_\lambda(x,w,t,s)=(x,\lambda w,\lambda t,\lambda s)$,
\[
 G\circ D_\lambda=|\lambda|^2G.
\]
Dilation therefore preserves critical points of restrictions to graphs, their real nondegeneracy and their Levi inertia, for every candidate metric. It maps the graph of a function $g$ over $C$ to the graph of $\lambda g\circ d_\lambda^{-1}$ over $d_\lambda(C)$, with $d_\lambda$ as in \cref{lem:dilated-countable-holomorphic-tests}. A function within $|\lambda|\delta_j$ of $\lambda g_j\circ d_\lambda^{-1}$ on $d_\lambda(C_j)$ pulls back under the dilation to a function within $\delta_j$ of $g_j$ on $C_j$. In particular, the metric neighborhood $\mathcal V_j$ needs no change.

Applying \cref{lem:dilated-countable-holomorphic-tests}, we obtain one function $f\in\mathcal O(B)$ that satisfies all the dilated tests. Its graph $\mathcal S$ is biholomorphic to the connected affine Stein manifold $B$, and it is smooth and closed in the open subset $\{x\in U,\ t\neq0\}$ of $E^\times$. Every candidate metric belongs to some $\mathcal V_j$, so its restriction to $\mathcal S$ has the required nondegenerate critical point. Since $\dim_{\mathbb C}\mathcal S=n+r-1$, at least $r$ positive Levi eigenvalues leave at most $n-1$ negative ones.
\end{proof}

\Needspace{10\baselineskip}
\begin{corollary}\label{cor:hirzebruch-fixed-stein-hypersurface}
\index{Hirzebruch surface!fixed Stein hypersurface}
For each $e\geq0$, the bundle $E$ of \cref{cor:hirzebruch-stable-ample-finsler-obstruction} can be chosen so that $E^\times$ contains a connected smooth Stein hypersurface on which the restriction of every smooth strongly pseudoconvex squared Finsler metric on $E$ has a nondegenerate critical point with at least two positive Levi eigenvalues and at most one negative Levi eigenvalue.
\end{corollary}
\begin{proof}
The bundle constructed in the proof of \cref{thm:hirzebruch-stable-ample-approximation} is a fiber of an algebraic family on a chart of the prioritary stack, so it is algebraic. Apply \cref{thm:fixed-stein-finsler-hypersurface} with $n=r=2$, using the curvature obstruction of \cref{cor:hirzebruch-stable-ample-finsler-obstruction}.
\end{proof}

% Source reconstruction: Liu--Wan 2609.37309v1, Proposition 5.1, pp. 24--27.
% Algebraicity is explicit here; the general holomorphic-to-algebraic GAGA step is not needed.
% Classical accepted inputs: Oka--Weil, one-variable Runge, affine Stein theory.
% The complement-convexity application and Fulton primary proof are not certified here.
